% META: {"id":"1SPE-PROBCOND-CO-016","chapitre":"1SPE-PROBA-COND","type_objet":"corrige","exercice_ref":"1SPE-PROBCOND-EX-016","capacites_codes":["C2"],"sources_inspiration":[],"status":"approved","origin":"generated","status_history":["generated","audited","accepted","approved"],"release_acceptance":"RELEASE_OWNER_BATCH_ACCEPTANCE","acceptance_closure_digest":"sha256:9b3ccf9a81c5520fbb7e03b7b2d3e7bf2057a02834b6d908bf3be5f49f3b553f","acceptance_receipt_digest":"sha256:4fad86f0282e0fa4e0419cca1ad6379ae7af5a280c04a4a90880f90a1c044242"}
% BEGIN-VERIFY
% from sympy import Rational
% P_meme = Rational(3, 7)
% P_aumoins1B = 1 - Rational(2, 7)
% assert P_aumoins1B == Rational(5, 7)
% END-VERIFY
\begin{corrige}{1SPE-PROBCOND-EX-016}

\textbf{1.} Arbre construit avec $P(R_1)=\frac{4}{7}$, $P(B_1)=\frac{3}{7}$ et les conditionnelles sans remise.

\medskip
\textbf{2.} $P(\text{meme}) = \frac{4}{7}\times\frac{3}{6} + \frac{3}{7}\times\frac{2}{6} = \frac{2}{7} + \frac{1}{7} = \frac{3}{7}$.

\medskip
\textbf{3.} $P(\text{au moins 1 bleue}) = 1 - P(R_1 \cap R_2) = 1 - \frac{2}{7} = \frac{5}{7}$.

\end{corrige}
